\subsection{LLMs in der Umfrageerstellung} \label{sec:llm}

\acfp{llm} \textit{(dt. große Sprachmodelle)} sind KI-Modelle, die auf großen Textmengen trainiert wurden und in der Lage sind, menschenähnliche Texte zu generieren. In den letzten Jahren wurden große Fortschritte in dem Feld des \ac{nlp} erzielt, was zu einem großen Anstieg der Nutzung von KI-Chatbots geführt hat. Insbesondere mit der Veröffentlichung von OpenAIs ChatGPT im November 2022 wurde ein neuer Meilenstein erreicht. \acp{llm} haben sich als äußerst leistungsfähig erwiesen und sind in der Lage, komplexe Aufgaben zu bewältigen, die zuvor menschliche Intelligenz erforderten. Diese Modelle können nicht nur Texte generieren, sondern auch Fragen beantworten, Texte zusammenfassen und sogar kreative Inhalte erstellen. \cite{llm-grundlagen}

Die Entwicklung von \acp{llm} kann in drei Phasen unterteilt werden \cite{llm-grundlagen}:
\begin{enumerate}
    \item \textbf{Training Phase:} In dieser Phase werden die Modelle auf großen vorbereiteten Textmengen trainiert, um die Muster und Strukturen der Sprache zu lernen. Das Training erfolgt in einem unüberwachtem Lernprozess basierend auf der Transformer Architektur, die 2017 von Vaswani et al. eingeführt wurde. Diese Architektur stützt sich auf einen "Attention Mechanism", um den Kontext von Wörtern in einem Satz zu verstehen und relevante Informationen zu extrahieren. \cite{llm-attention} Dabei werden die erlernten Informationen numerisch in den Parametern des Modells gespeichert. Die Größe der Modelle wird in der Regel in Milliarden von Parametern angegeben, was auf die Komplexität und den Umfang des Modells hinweist. \cite{llm-grundlagen}
    \item \textbf{Fine-Tuning:} In dieser Phase wird das vortrainierte Modell auf spezifische Aufgaben oder Domänen angepasst. Dies geschieht häufig durch überwachtes Lernen, bei dem das Modell mit gekennzeichneten Daten trainiert wird. Weitere bekannte Methoden sind Reinforcement Learning from Human Feedback (RLHF), Transfer Learning und Feedback Loops, wo durch Nutzereingaben und Rückmeldungen das Modell kontinuierlich verbessert wird. \cite{llm-privacy} 
    \item \textbf{Inferenz:} In dieser Phase wird das trainierte Modell verwendet, um Vorhersagen oder Texte zu generieren. Dies geschieht in der Regel in Echtzeit und erfordert eine effiziente Hardware-Infrastruktur.
\end{enumerate} 

Da \acp{llm} in der Lage sind, menschenähnliche Sprache zu generieren und komplexe Zusammenhänge zu verstehen, eignen sie sich hervorragend für die Unterstützung bei der Umfrageerstellung. Mit dem entsprechenden Hintergrundwissen über die Zielgruppe sind sie in der Lage, präzise und relevante Fragen für die Nutzenden zu formulieren. Mit dieser Funktion können sie aus einem trivialem Prompt den Kontext einer Umfrage analysieren und passende Fragen erstellen. Durch die Automatisierung wird der Prozess der Umfrageerstellung erheblich beschleunigt und vereinfacht. Neben dem genannten Vorteil besteht zudem die Möglichkeit, dass \acp{llm} qualitativ hochwertigere Fragen generieren können. Durch ihr Vorwissen über die objektiven Richtlinien zur Umfragenerstellung sind sie in der Theorie in der Lage, präzisere und weniger verwirrende Fragen mit passenden Antworttypen zu formulieren.  

\input{\currfiledir 01_HW_Anforderungen.tex}
\input{\currfiledir 02_Datenschutz.tex}
\input{\currfiledir 03_Prompt_Engineering.tex}